\documentclass[11pt]{article}

% ====================================================
% PACKAGES
% ====================================================
\usepackage[margin=1in, headheight=13.6pt]{geometry}
\usepackage[T1]{fontenc}
\usepackage{lmodern}
\usepackage{microtype}
\usepackage{setspace}
\usepackage{amsmath, amssymb, mathtools}
\usepackage{siunitx}
\usepackage{graphicx}
\usepackage{float}
\usepackage{enumitem}
\usepackage{booktabs}
\usepackage{array}
\usepackage{xcolor}
\usepackage{hyperref}
\usepackage{fancyhdr}
\usepackage{lastpage}
\usepackage{indentfirst}
\usepackage[justification=centering]{caption}

\pagestyle{fancy}
\fancyhf{}
\lhead{\Assignment\ |\ \course}
\rhead{\Name\ |\ \today}
\cfoot{Page \thepage\ of~\pageref{LastPage}}

\hypersetup{
    colorlinks,
    linkcolor={red!50!black},
    citecolor={blue!50!black},
    urlcolor={blue!80!black}
}

\newcommand{\Name}{Arael Anaya}
\newcommand{\Assignment}{2.2 Payload Driven Mission Design}
\newcommand{\course}{Space Robotics}

% ====================================================
% DOCUMENT
% ====================================================
\begin{document}

% ============================================================
\noindent\textbf{Payload: REASON (Radar for Europa Assessment and Sounding: Ocean to Near-surface)}
% ============================================================

\medskip

REASON is a dual-frequency ice-penetrating radar sounder currently flying aboard NASA's Europa
Clipper spacecraft. The instrument was developed by researchers at the University of Texas
Institute for Geophysics and fabricated by JPL, and it launched on October 14, 2024 atop a
SpaceX Falcon Heavy rocket bound for Europa~[3]. REASON operates at two frequency bands
simultaneously, allowing it to distinguish between near-surface structure and deep basal
reflections, a capability that no previous Europa-targeted instrument has demonstrated~[1].
The instrument is flight-proven hardware with full mission-level heritage, placing it at TRL~9.

Europa Clipper carries REASON as part of a flyby science campaign: the spacecraft will loop
past Europa roughly 50 times, collecting wide-swath regional maps of the ice shell from
altitude~[2]. That is exactly what Clipper is designed to do, and it will do it well. What
Clipper is not doing is surface-level, high-resolution traverse profiling targeted at a specific
candidate cryobot insertion site. There is no rover, no in-situ localization, and no onboard
site-selection logic. The orbital geometry also limits any single flyby to a brief snapshot
rather than a continuous along-track profile. My mission picks up where Clipper leaves off: a
surface rover carrying an adapted REASON-heritage sounder drives a continuous traverse, building
a spatially registered subsurface map at 100~m resolution or better, and autonomously flags the
thinnest viable ice for a future cryobot deployment.

\bigskip

\begin{enumerate}[label=\Roman*., leftmargin=2.5em, itemsep=0.75em]

    \item \textbf{What parameters are measured?}

    REASON measures four primary parameterse. \textbf{Two-way travel time}~(ns)
    is the elapsed time between the transmitted pulse and the received basal reflection. The
    raw measurement from which ice thickness is derived. \textbf{Signal amplitude at the basal
    reflector}~(dB) captures the strength of the return from the ice-ocean interface. Returns
    too weak to distinguish from the noise floor are flagged and excluded from site-selection
    logic. \textbf{Bulk ice dielectric permittivity} (dimensionless) is estimated from the
    travel time and signal characteristics, where pure water ice sits around 3.15 and deviations
    indicate brine inclusions or heterogeneous layering. Finally, \textbf{derived ice shell
    thickness}~(m) is computed onboard as $d = c \cdot t \,/\, 2\sqrt{\varepsilon_r}$, where
    $c$ is the speed of light, $t$ is the two-way travel time in seconds, and $\varepsilon_r$
    is the estimated permittivity. This is the primary science output that feeds the subsurface
    map and the site-selection filter.

    \item \textbf{What would a single data entry or record look like?}

    The data product is organized as a multi-sheet spreadsheet with one row per radar trace,
    time-aligned across sheets by \texttt{trace\_id}. Each record captures the full measurement
    context: radar sounder outputs, rover position, terrain conditions, and system health. A
    data quality flag (NOMINAL / CLUTTER / NOISE\_FLOOR / INVALID) is attached to every trace
    so the onboard planner and downstream scientists can immediately distinguish usable science
    data from degraded returns.

    \smallskip
    \noindent\textit{Example single record (Radar\_Traces sheet):}

    \begin{table}[H]
    \centering
    \small
    \begin{tabular}{ll}
    \toprule
    \textbf{Field} & \textbf{Value} \\
    \midrule
    \texttt{trace\_id}         & TR-0001 \\
    \texttt{sol}               & 1 \\
    \texttt{timestamp\_s}      & 94.4~s \\
    \texttt{pos\_x\_m}         & 93.8~m (local ENU east) \\
    \texttt{pos\_y\_m}         & $-7.1$~m (local ENU north) \\
    \texttt{along\_track\_m}   & 94.0~m \\
    \texttt{twt\_ns}           & 21647.3~ns \\
    \texttt{amplitude\_dB}     & $-41.62$~dB \\
    \texttt{permittivity}      & 3.0624 \\
    \texttt{ice\_thickness\_m} & 1855.5~m \\
    \texttt{tid\_krad}         & 0.542~krad \\
    \texttt{quality\_flag}     & NOMINAL \\
    \bottomrule
    \end{tabular}
    \end{table}

    \item \textbf{Who is the intended user of this data?}

    The primary real-time consumer is the \textbf{onboard autonomous planning algorithm}.
    After each consensus round, the planner ingests the daily summary sheet to decide whether
    the traverse has produced a viable candidate site (ice thickness ${<}$5~km, slope
    ${<}20^{\circ}$, roughness ${<}$0.7). If the criteria are met, it triggers a cryobot
    deployment recommendation; if not, it generates the next traverse waypoint. No human is in
    this loop. The secondary long-term consumer is \textbf{planetary scientists}: the full
    four-sheet data product, downlinked to Earth over the mission lifetime, provides a spatially
    registered along-track profile of ice shell thickness and permittivity variation across
    Europa's surface. This feeds directly into habitability assessments and informs the design
    of future subsurface access missions, extending the science value of the dataset well beyond
    the immediate site-selection objective.

\end{enumerate}

\bigskip

\noindent\textbf{\underline{References:}}

\begin{enumerate}[label={[\arabic*]}, leftmargin=3em, itemsep=0.25em]

    \item NASA Jet Propulsion Laboratory, ``Europa Clipper,'' JPL, Pasadena, CA.
    [Online]. Available: \url{https://www.jpl.nasa.gov/missions/europa-clipper/}.
    [Accessed: Jun.\ 16, 2026].

    \item Wikipedia contributors, ``Europa Clipper,'' \textit{Wikipedia, The Free Encyclopedia}.
    [Online]. Available: \url{https://en.wikipedia.org/wiki/Europa_Clipper}.
    [Accessed: Jun.\ 16, 2026].

    \item University of Texas at Austin Jackson School of Geosciences,
    ``Scientific Instrument Developed by UTIG Scientists Headed to Europa,''
    UT Austin, Austin, TX, Apr.\ 2025. [Online]. Available:
    \url{https://www.jsg.utexas.edu/news/2025/04/on-oct-14-2024-nasas-europa-clipper-successfully-launched-atop-a-spacex-falcon-heavy-rocket-bound-for-jupiters-moon-europa/}.
    [Accessed: Jun.\ 16, 2026].

\end{enumerate}

\end{document}